\begin{table}[!ht]
\centering
\begin{threeparttable}
\caption{Directional Accuracy versus Naive Benchmark}
\label{tab:naive_comparison}

\begin{tabular}{llrrlr}
\toprule
Model & Horizon & DA & Naive DA & Beats Naive & DA Difference\\
\midrule
OLS & 1d & 0.3219048 & 0.2761905 & Yes & 0.0457\\
OLS & 3d & 0.4761905 & 0.3200000 & Yes & 0.1562\\
OLS & 5d & 0.4990476 & 0.3028571 & Yes & 0.1962\\
Random Forest & 1d & 0.2895200 & 0.2761900 & Yes & 0.0133\\
Random Forest & 3d & 0.4666700 & 0.3200000 & Yes & 0.1467\\
Random Forest & 5d & 0.4838100 & 0.3028600 & Yes & 0.1810\\
XGBoost & 1d & 0.2781000 & 0.2761900 & Yes & 0.0019\\
XGBoost & 3d & 0.4533300 & 0.3200000 & Yes & 0.1333\\
XGBoost & 5d & 0.4819000 & 0.3028600 & Yes & 0.1790\\
SVM & 1d & 0.3200000 & 0.2761900 & Yes & 0.0438\\
SVM & 3d & 0.4800000 & 0.3200000 & Yes & 0.1600\\
SVM & 5d & 0.4819000 & 0.3028600 & Yes & 0.1790\\
ANN & 1d & 0.3466700 & 0.2761900 & Yes & 0.0705\\
ANN & 3d & 0.4685700 & 0.3200000 & Yes & 0.1486\\
ANN & 5d & 0.4819000 & 0.3028600 & Yes & 0.1790\\
\bottomrule
\end{tabular}
\begin{tablenotes}
\footnotesize
\item \textit{Note.} DA Difference = Model DA minus Naive DA. Positive values indicate the model outperforms the naive benchmark. The naive benchmark uses the previous day BPI log return (lag1) as the prediction.
\end{tablenotes}
\end{threeparttable}
\end{table}
